\subsection{FREE MAGMAS}

A sequence of sets \(\mathbf{M}_{n}(\mathbf{X})\) is defined by induction on the integer \(n \geqslant 1\) as follows: writing \(\mathbf{M}_{\mathbf{1}}(\mathbf{X}) = \mathbf{X}\) , for \(n \geqslant 2\) , \(\mathbf{M}_{n}(\mathbf{X})\) is the set the sum of the sets \(\mathbf{M}_{p}(\mathbf{X}) \times \mathbf{M}_{n-p}(\mathbf{X})\) for \(1 \leqslant p \leqslant n - 1\) . The set the sum of the family \((\mathbf{M}_{n}(\mathbf{X}))_{n \geqslant 1}\) is denoted by \(\mathbf{M}(\mathbf{X})\) ; each of the sets \(\mathbf{M}_{n}(\mathbf{X})\) is identified with its canonical image in \(\mathbf{M}(\mathbf{X})\) . For every element w of \(\mathbf{M}(\mathbf{X})\) there exists a unique integer n such that \(w \in \mathbf{M}_{n}(\mathbf{X})\) ; it is called the length of w and denoted by \(l(w)\) . The set X consists of the elements in \(\mathbf{M}(\mathbf{X})\) of length 1.

Let \(w\) and \(w'\) be in \(\mathbf{M}(\mathbf{X})\) ; write \(p = l(w)\) and \(q = l(w')\) . The image of \((w, w')\) under the canonical injection of \(\mathbf{M}_p(\mathbf{X}) \times \mathbf{M}_q(\mathbf{X})\) into the sum set \(\mathbf{M}_{p + q}(\mathbf{X})\) is called the composition of \(w\) and \(w'\) and is denoted by \(ww'\) or \(w.w'\) . Then \(l(w.w') = l(w) + l(w')\) and every element of \(\mathbf{M}(\mathbf{X})\) of length \(\geqslant 2\) can be written uniquely in the form \(w'w''\) with \(w', w''\) in \(\mathbf{M}(\mathbf{X})\) .

The set \(\mathbf{M}(\mathbf{X})\) with the law of composition \((w, w') \mapsto w.w'\) is called the free magma constructed on \(\mathbf{X}\) (§ 1, no. 1, Definition 1).

PROPOSITION 1. Let M be a magma. Every mapping f of X into M may be extended in a unique way to a morphism of M(X) into M.

By induction on \(n \geqslant\) , mappings \(f_n: \mathbf{M}_n(\mathbf{X}) \to \mathbf{M}\) are defined as follows: let \(f_{1} = f\) ; for \(n \geqslant 2\) , the mapping \(f_{n}\) is defined by \(f_{n}(w, w') = f_{p}(w).f_{n-p}(w')\) for \(p = 1, 2, \ldots, n - 1\) and \((w, w')\) in \(\mathbf{M}_{p}(\mathbf{X}) \times \mathbf{M}_{n-p}(\mathbf{X})\) . Let \(g\) be the mapping of \(\mathbf{M}(\mathbf{X})\) into \(\mathbf{M}\) which induces \(f_{n}\) on \(\mathbf{M}_{n}(\mathbf{X})\) for every integer \(n \geqslant 1\) . Clearly \(g\) is the unique morphism of \(\mathbf{M}(\mathbf{X})\) into \(\mathbf{M}\) which extends \(f\) .

Let \(u\) be a mapping of \(\mathbf{X}\) into a set \(\mathbf{Y}\) . By Proposition 1, there exists one and only one homomorphism of \(\mathbf{M}(\mathbf{X})\) into \(\mathbf{M}(\mathbf{Y})\) which coincides with \(u\) on \(\mathbf{X}\) . It will be denoted by \(\mathbf{M}(u)\) . If \(v\) is a mapping of \(\mathbf{Y}\) into a set \(\mathbf{Z}\) , the homomorphism \(\mathbf{M}(v) \circ \mathbf{M}(u)\) of \(\mathbf{M}(\mathbf{X})\) into \(\mathbf{M}(\mathbf{Z})\) coincides with \(v \circ u\) on \(\mathbf{X}\) , whence

\[
\mathbf {M} (v) \circ \mathbf {M} (u) = \mathbf {M} (v \circ u).
\]

PROPOSITION 2. Let \(u: \mathbf{X} \to \mathbf{Y}\) be a mapping. If \(u\) is injective (resp. surjective, bijective), so is \(\mathbf{M}(u)\) .

Suppose \(u\) is injective. When \(X\) is empty, \(M(X)\) is empty, hence \(M(u)\) is injective. If \(X\) is non-empty, there exists a mapping \(v\) of \(Y\) into \(X\) such that \(v \circ u\) is the identity mapping of \(X\) (Set Theory, II, § 3, no. 8, Proposition 8); the mapping \(M(v) \circ M(u) = M(v \circ u)\) is the identity mapping of \(M(X)\) and hence \(M(u)\) is injective.

When \(u\) is surjective, there exists a mapping \(w\) of \(Y\) into \(X\) such that \(u \circ w\) is the identity mapping of \(Y\) (Set Theory, II, §3, no. 8, Proposition 8). Then \(M(u) \circ M(w) = M(u \circ w)\) is the identity mapping of \(M(Y)\) and hence \(M(u)\) is surjective.

Finally, if \(u\) is bijective, it is injective and surjective and hence \(\mathbf{M}(u)\) has the same properties.

Let S be a subset of X. By Proposition 2 the injection of S into X can be extended to an isomorphism of M(S) onto a submagma M'(S) of M(X). The magmas M(S) and M'(S) are identified by means of this isomorphism. Then M(S) is the submagma of M(X) generated by S.

Let X be a set and \((u_{\alpha}, v_{\alpha})_{\alpha \in I}\) be a family of ordered pairs of elements of M(X). Let R be the equivalence relation on M(X) compatible with the law of M(X) and generated by the \((u_{\alpha}, v_{\alpha})\) (§1, no. 6). The magma M(X)/R is called the magma defined by X and the relators \((u_{\alpha}, v_{\alpha})_{\alpha \in I}\) . Let h be the canonical morphism of M(X) onto M(X)/R. Then M(X)/R is generated by \(h(X)\) .

Let N be a magma and \((n_{x})_{x\in\mathbf{X}}\) a family of elements of N. Let k be the morphism from \(\mathbf{M}(\mathbf{X})\) to N such that \(k(x)=n_{x}\) for all \(x\in X\) (Proposition 1). If \(k(u_{\alpha})=k(v_{\alpha})\) for all \(\alpha\in I\) , there exists one and only one morphism \(f\colon\mathbf{M}(\mathbf{X})/\mathbf{R}\to\mathbf{N}\) such that \(f(h(x))=n_{x}\) for all \(x\in X\) (§ 1, no. 6, Proposition 9).

